\section{从“写完成”到“掉电后仍存在”}

写系统调用返回、块请求完成、设备缓存接收、介质完成编程是四个不同边界。若应用先写数据再写指针，掉电后不能允许“新指针已持久化而新数据尚未持久化”。日志、写时复制和显式 flush 都是在安排这两个持久化顺序。

\begin{longtable}{L{0.2\linewidth}L{0.34\linewidth}L{0.36\linewidth}}
\toprule
层次 & 完成通常表示 & 尚未保证 \\
\midrule
页缓存 & 数据已被内核接收 & 块层、设备和介质完成 \\
块层完成 & 设备报告请求完成 & 无断电保护缓存中的数据抗掉电 \\
设备 FLUSH 完成 & 之前写入跨过设备规定的持久化边界 & 应用自己的元数据顺序必然正确 \\
应用事务提交 & 日志/数据顺序满足协议 & 未纳入协议的外部系统状态 \\
\bottomrule
\end{longtable}

\section{校验、冗余与故障域}

ECC 修复单个器件或扇区中的有限错误，校验和发现数据路径或静默损坏，镜像与纠删码在设备失效后重建数据；它们解决的故障粒度不同。RAID 不是备份：控制器错误、软件误删、勒索软件和整机电源事故可能同时影响阵列所有成员。

分析可靠性时先画故障域：同一芯片、同一通道、同一背板、同一电源、同一机箱、同一机房。两份数据若共享同一个失效点，就不能简单按两次独立概率相乘。重建期间还会增加介质流量与暴露时间，因此容量、校验强度、备用空间和重建限速必须一起设计。
